\documentclass{usstthesis}

\usstthesisset{%
    session=2015,   % 2015 届
    title=上海理工大学本科毕业设计(论文) \LaTeX 文档类, % 论文中文标题
    titleen=A \LaTeX{} document class for the Final Project/Thesis for Undergraduates of USST,  % 论文英文标题
    institute=光电信息与计算机工程学院, % 学院
    major=计算机科学与技术, % 专业
    name=你的名字,  % 姓名
    number=1234567, % 学号
    mentor=指导教师,    % 指导教师
    date=2019年1月1日   % 完成日期
}

\lstset{language=C}

\renewcommand{\arraystretch}{1.1}   % 表格行高 1.1 倍

\begin{document}
\outputfrontmatter  % 输出封面、承诺书、摘要和目录
% 注意，摘要请在 ./elements/abstract.tex 编辑
    
\mainmatter % 中间内容
% 可以使用分文件编辑，最后 input 每个章节
\chapter{论文的内容要求}
\par 如要使用宋体粗体，请使用 \verb|{\bSong 要加粗的字}|：
{\bSong 一份完整的毕业设计（论文）包括：标题、基本信息、承诺书、摘要、关键词、目录、正文、参考文献、致谢、附录等。}

\section{标题}
\par 标题包括中文标题和外文标题。

\subsection{中文标题}
\par 中文标题应该简短、明确、有概括性，不宜超过20个汉字。

\subsection{外文标题}
\par 外文标题应该简短、明确、翻译正确。

\section{基本信息}
\par 基本信息包括：学院名称、专业名称、作者姓名、作者学号、指导教师姓名及职称，以及论文完成的日期。

\section{承诺书}
\par 承诺书是论文作者对学术诚信的庄重承诺。本规范提供了上海理工大学本科毕业设计(论文)承诺书的一个规范文本，作者在认真仔细阅读后签上姓名和日期。

\section{摘要}
\par 摘要包括中文摘要和外文摘要两部分。
\par 中外文摘要均包括摘要正文和关键词。

\subsection{摘要正文}
\par 论文摘要简要陈述本科毕业设计(论文)的内容，创新见解和主要论点。中文摘要在500字左右，外文摘要应与中文摘要的内容相符。

\subsection{关键词}
\par 关键词是反映毕业设计(论文)主题内容的名词，是供检索使用的。关键词条应为通用词汇，不得自造关键词。关键词一般为3至5个，按其外延层次，由高至低顺序排列。关键词排在摘要正文部分下方。

\section{目录}
\par 目录按三级标题编写，要求标题层次清晰，并标明页码。

\section{正文}
\par 正文篇幅要求15000字以上（其中，英语、德语专业毕业设计（论文）应不少于5000外文单词，日语专业应不少于8000日语假名，艺术设计类专业不少于3000字。毕业设计（论文）的核心设计、研究篇幅应占全篇幅的三分之二以上。）。内容包括绪论、正文主体与结论。

\subsection{绪论}
\par 绪论是研究工作的概述，内容包括：本课题的意义、目的、研究范围及要达到的技术要求；简述本课题在国内外的发展概况及存在的问题。
\par 绪论一般作为在毕业设计（论文）正文的第1章，并在一章内完成。

\subsection{正文主体}
\par 正文主体是研究工作的详述，内容包括：问题的提出，研究工作的基本前提、假设和条件；模型的建立，实验方案的拟定；基本概念和理论基础；设计计算的主要方法和内容；实验方法、内容及其分析；理论论证及其应用，研究结果，以及对结果的讨论等。
\par 正文主体一般可分为若干章完成。

\subsection{结论}
\par 结论是研究工作的总结，内容包括：对所得结果与已有结果的比较和课题尚存在的问题，以及进一步展开研究的见解与建议。
\par 结论一般作为论文正文的最后一章，并在一章内完成。

\section{参考文献}
\par 参考文献为研究中参考的资料，包括专著、论文、年鉴、网站等。所引用的文献必须是公开发表的与毕业设计（论文）工作直接有关的文献，且经过本人阅读理解。列入的主要文献要求不少于6篇，其中外文文献不少于2篇。

\section{致谢}
\par 毕业设计（论文）是在指导教师的指导下完成，理应致谢。还应对完成论文提供过帮助的其他人员致谢。切忌泛滥和溢美。

\section{附录}
\par 附录是对于一些不宜放在正文中，但又直接反映研究工作的材料（如设计图纸、实验数据、计算机程序等）附于文本末尾。
 % 论文的内容要求
% 也可以直接在此文件中编写，但是比较乱
\chapter{论文的格式要求}
\par 论文的格式要求包括：纸张大小、纸张方向、页边距、版式、文档网格、字体与字号、段落和间距等\cite{机器学习}。
\par 建议采用 \LaTeX2e 语法和 \verb|xelatex| 引擎编译。
\par 由于论文格式问题非常繁杂，无法将所有设置描述清楚，只能对一些主要的设置做出扼要的说明。一个快捷有效的方法就是把本规范的电子版作为模板。
\par 小贴士：使用 \verb|xelatex| 编译时，可以直接生成双面打印的 PDF 文件。如果需要编译两次生成交叉引用，第一次编译时加入 \verb|--no-pdf| 选项可以加快编译速度。

\lstinputlisting[caption=code.c]{code/code.c}

\captionof{lstlisting}{main.h}
\begin{lstlisting}
int main(int argc, char* argv[]);
\end{lstlisting}

\section{页面设置}
\subsection{纸张}
\par 纸张大小：A4。
\par 纸张方向：纵向。

\subsection{页边距}
\par 页边距；上2.5厘米，下2.5厘米，内侧3厘米，外侧2.5厘米。
\par 页码范围：对称页边距。

\subsection{版式}
\par 节：奇数页。
\par 页眉和页脚：奇偶页不同，距边界：页眉1.5厘米，页脚1.75厘米。

\subsection{文档网格}
\par 网格：无网格。

\subsection{字体}
\par 中文字体：宋体，使用{\bSong 宋体粗体}时，请使用 \verb|{\bSong 要加粗的文字}|。
\par 西文字体：Times New Roman。
\par 字形：常规。
\par 字号：小四。

\subsection{段落}
\par 对齐方式：两端对齐。
\par 首行缩进：2字符。
\par 行距：多倍行距1.25。

\section{封面}
\subsection{标题}
\par 中文标题：二号华文中宋和 Times New Roman 加粗，居中，左、右侧缩进均为4字符。
\par 外文标题：小二号 Times New Roman 加粗，居中，左、右侧缩进均为4字符。

\subsection{基本信息}
\par 基本信息是一个表格，左列为基本信息名称，右列为需要填写的基本信息。
\par 基本信息：四号华文中宋和 Times New Roman 加粗居中。

\section{承诺书}
\par 承诺书：三号华文中宋加粗，居中，段前4行，段后2行。
\par 承诺书文本：小四号宋体和 Times New Roman，首行缩进2字符，1.25倍行距。

\section{摘要}
\par 摘要：三号华文中宋加粗，居中，段前4行，段后2行。
\par 摘要文本：小四号宋体和 Times New Roman，首行缩进2字符，1.25倍行距。
\par 摘要文本结束后空一行。
\par 关键词：小四号宋体加粗顶格，××××：小四号宋体和 Times New Roman，各关键词之间2空格。

\section{ABSTRACT}
\par ABSTRACT：三号Times New Roman加粗，居中，段前4行，段后2行。
\par ABSTRACT文本：小四号Times New Roman，首行缩进2字符，1.25倍行距。
\par ABSTRACT文本结束后空一行。
\par KEY WORDS：小四号Times New Roman加粗顶格，××××：小四号Times New Roman，各关键词之间2空格。

\section{目录}
\par 目录：三号华文中宋加粗，居中，段前4行，段后2行。
\par 以下内容用小四号宋体和 Times New Roman，1.25倍行距：
\par 摘要：加粗，首行缩进2字符，段前0.5行；
\par ABSTRACT：加粗，首行缩进2字符，段前0.5行；
\par 第1章 ××××.........................1：加粗，首行缩进2字符，段前0.5行；
\par \quad 1.1 ××××............................1：加粗，首行缩进3字符；
\par \qquad 1.1.1 ××××.....................1：加粗，首行缩进4字符。

\section{正文}
\par 一级标题：三号华文中宋和 Times New Roman 加粗，居中，段前4行，段后2行。
\par 二级标题：四号宋体和 Times New Roman 加粗，左对齐顶格，段前0.5行，段后0行。
\par 三级标题：小四号宋体和 Times New Roman 加粗，左对齐顶格，段前0.5行，段后0行。
\par 正文文字：小四号宋体和 Times New Roman，首行缩进2字符，1.25倍行距。

\section{参考文献}
\par 参考文献：三号华文中宋加粗，居中，段前4行，段后2行。
\par 参考文献序号用方括号括起。
\par 参考文献序号和内容用五号宋体和 Times New Roman。

\section{致谢}
\par 致谢：三号华文中宋加粗，居中，段前4行，段后2行。
\par 致谢文本：小四号宋体和 Times New Roman，首行缩进2字符，1.25倍行距。
 % 论文的格式要求
\chapter{论文的写作细则}
\section{书写}
\par 正文中的任何部分，如表、图，应限制在版心的范围以内，不要超出版心的范围。
\par 汉字必须使用国家公布的规范字。

\section{标点符号}
\par 标点符号按新闻出版署公布的“标点符号用法”使用。

\section{名词、名称}
\par 科学技术名词术语采用全国自然科学名词审定委员会公布的规范词或国家、部颁标准中规定的名称，尚未统一规定的名词术语，可采用惯用的名称。使用外文缩写代替某一名词术语时，在首次出现处加括号注明其含义。外国人名一般用英文原名，按名前姓后的原则书写。一般很熟知的外国人名（如牛顿、达尔文、马克思等）可按通常标准译法书写中文译名。

\section{量和单位}
\par 量和单位采用中华人民共和国的国家标准GB3100～GB3102-93。非物理量的单位（如件、台、人、元等），可用汉字与符号构成组合形式的单位，如：件/台、元/km等。

\section{数字}
\par 测量统计数据一律用阿拉伯数字。行文叙述个位数时，可用中文数字，如“他发现两颗小行星”、“三力作用于一点”，不宜写成“他发现2颗小行星”、“3力作用于1点”。约数可用中文数字，也可以用阿拉伯数字，如“约一百二十八人"，也可写成"约128人”。

\section{标题层次}
\par 标题层次采用采用三级数字编号方法，例如第一级为“第1章”，第二级为“1.1”、“1.2”，第三级为“1.1.1”、“1.1.2”等。层次控制在三级以内，必要时可增设四级。两级之间用下角圆点隔开，每级末尾不加标点。
\par 各层标题均单独占行书写。第一级（章）标题位于新起始页上方正中，第二、三级其标题顶格书写，后空一格书写标题，末尾不加标点。
正文中对总项包括的分项采用(1)、(2)、(3)…的序号，对分项中的小项采用1)、2)、3)...的序号，序号后不再加其他标点，序号前空二格书写。

\section{注释（脚注）}
\par 个别名词或情况需要解释时，可加注说明，注释可用页末注（注文放在加注页的下端）或篇末注（全部注文集中在正文末尾），而不可用行中注（注文夹在正文中间）。

\section{公式}
\par 公式居中书写，建议使用 \verb|equation| 等公式环境，公式较长时应在“$=$”前转行或在“$+$、$-$、$\times$、$\div$”等运算符号处转行，等号或运算符号应在转行后的行首。公式的编号用圆括号括起来放在公式右边行末，公式和编号之间留空。微分算子可使用 \verb|usstthesis| 文档类定义的 \verb|\dif|。例：
\begin{equation}
    F(j\omega)=\int_{-\infty}^{\infty}f(t)e^{-j\omega t}\dif t
\end{equation}

\section{表格}
\par 每个表格应有表序和表题，表序和表题应在表格的上方，居中，使用表格环境的 \verb|\caption| 书写，便于添加引用，表格整体置于 \verb|table| 环境中。一般情况下表格采用五号宋体。例表 \ref{table:tsp calc}：
\begin{table}[htb]
    \tableCapSet    % 使用此命令调整 caption 间距
    \caption{TSP 问题的计算量（计算速度：$10^8/$s）}
    \label{table:tsp calc}
    \centering
    \zihao{5}
    \begin{tabular}{c|c|c}
        \hlineB{3}  % 线宽为3倍的横线
        城市数目 $n$ & 计算量               & 计算时间                \\
        \hlineB{2}  % 线宽为2倍的横线
        10           & $1.8 \times 10^5$    & $1.8 \times 10^{-3}$ 秒 \\
        \hline
        15           & $4.4 \times 10^{10}$ & $7.3$ 小时              \\
        \hline
        20           & $6.0 \times 10^{16}$ & $19$ 年                 \\
        \hline
        25           & $3.1 \times 10^{23}$ & $1.0 \times 10^7$ 世纪  \\
        \hlineB{3}
    \end{tabular}
\end{table}

\section{插图}
\par 每幅插图应有图序和图题，图序和图题放在图下方居中处，建议使用 \verb|figure| 环境，并且使用 \verb|\includegraphics| 插图。一般情况下插图采用五号宋体。例图 \ref{figure:model}：

\begin{figure}[htb]
    \figureCapSet
    \centering
    \includegraphics[width=.8\linewidth]{resource/egf.png}
    \caption{信息传播模型}
    \label{figure:model}
\end{figure}

\section{参考文献}
\par 参考文献是毕业设计（论文）中引用文献出处的目录表，一律放在文后。书写格式按国家标准GB7714-87规定。
\par 参考文献按正文中出现的先后统一用阿拉伯数字进行自然编号，序码用方括号括起。且在正文引用处最后一个字的右上角，用方括号标明此序号（如×××\cite{论文}，以便查找）。一篇论著在论文中多处被引用时，在参考文献目录表中只应出现一次，序号以第一次出现的位置为准。
\par 具体参考文献的标注格式可参考正规出版的论文杂志（集）。多数论文数据库可以直接生成 bibtex 引用格式。

\section{页面}
\par 从第1章起，奇数页页眉写毕业设计（论文）的题目，偶数页页眉写“上海理工大学本科毕业设计（论文）”，字体为小五号宋体。

\section{页码}
\par 从目录首页到目录末页，在页面底端外侧加注页码，页码为小五号 Times New Roman 小写罗马数字，即 i、ii、iii 等。
\par 从第1章正文起，在页面底端外侧加注页码，页码为小五号 Times New Roman 阿拉伯数字，即1、2、3等。

 % 论文的写作细则
% 表格、插图等内容在 elements/chapter3.tex 中有所体现

\backmatter % 结尾内容

\begin{citelist}
    \bibitem{机器学习}刘琴.机器学习[J].武汉工程职业技术学院学报,2001,13(2):41-44.
    \bibitem{论文}中国科学技术信息研究所.2018年度中国科技论文统计与分析[J].科学,2018,70(6):57-59.
\end{citelist}
\chapter{致\quad{}谢}
\par 本论文是在导师的悉心指导下完成的，本文作者在此谨表示衷心的感谢。
\par ××× 老师也对本论文给予了许多宝贵的意见和建议，在此表示深深的谢意。 % 致谢


\end{document}